\chapter{Referencial teórico}\label{cap:referencial-teorico}

Este capítulo apresenta os fundamentos conceituais e tecnológicos que sustentam o desenvolvimento do S.R.I. São tratados o processo de desenvolvimento adotado, as linguagens de modelagem utilizadas na documentação, as tecnologias empregadas na construção do aplicativo móvel e os conceitos de roteirização que orientam a funcionalidade central do sistema.

\section{Engenharia de software e o processo RUP}

A engenharia de software compreende o conjunto de métodos, técnicas e ferramentas que permitem produzir software de qualidade de forma sistemática, tratando o desenvolvimento como um processo organizado e não como uma atividade artesanal \cite{sommerville2019}. Adotar um processo definido é o que torna possível planejar as atividades, estimar esforços, controlar o andamento do projeto e garantir que os artefatos produzidos possam ser verificados por terceiros \cite{pressman2016}.

Entre os processos disponíveis, o Rational Unified Process (RUP) propõe uma abordagem iterativa e incremental, organizada em quatro fases sequenciais: concepção, elaboração, construção e transição. Cada fase é composta por iterações que produzem uma versão executável do sistema, o que permite validar os requisitos junto aos envolvidos antes de avançar para etapas de maior custo \cite{kruchten2003}. O processo é dirigido por casos de uso, o que significa que as funcionalidades do sistema são descritas sob a perspectiva de quem as utiliza e servem de fio condutor para todas as demais atividades de modelagem, de implementação e de teste.

O RUP organiza as atividades de desenvolvimento em disciplinas, entre as quais se destacam a modelagem de negócios, o levantamento e a gestão de requisitos, a análise e o projeto, a implementação, os testes e a implantação. Essa separação é particularmente adequada a projetos conduzidos por uma única pessoa, como é o caso do estágio descrito neste relatório, pois torna explícito o que deve ser produzido em cada etapa e evita que a atividade de programação absorva integralmente o tempo disponível \cite{kruchten2003}. A documentação produzida nas disciplinas iniciais, formalizada nos artefatos de visão, de pedido do investidor, de especificação complementar e de especificação de casos de uso, estabelece a base a partir da qual o sistema é construído \cite{sriEspecComp2025}.

\section{UML e a modelagem do sistema}

A Unified Modeling Language (UML) é uma linguagem de modelagem padronizada que oferece notações gráficas para representar tanto a estrutura estática quanto o comportamento dinâmico de um sistema orientado a objetos \cite{omg2017uml}. Seu propósito é fornecer um vocabulário comum entre analistas, desenvolvedores e demais envolvidos, de modo que a descrição do sistema possa ser compreendida independentemente das decisões de implementação \cite{guedes2011}.

Entre os diagramas previstos pela linguagem, o diagrama de casos de uso representa as funcionalidades do sistema e as interações entre os atores externos e essas funcionalidades, sendo o artefato mais diretamente associado ao levantamento de requisitos \cite{guedes2011}. O diagrama de classes descreve as classes do domínio, seus atributos, seus métodos e os relacionamentos de associação, agregação e composição existentes entre elas, servindo de base para a implementação e para o projeto do banco de dados. O diagrama de sequência representa a troca de mensagens entre objetos ao longo do tempo, o que permite verificar a ordem das operações necessárias para realizar um caso de uso. O diagrama de estados descreve os estados possíveis de um objeto e as transições provocadas por eventos, sendo útil para modelar o ciclo de vida de entidades que mudam de situação durante a execução. O diagrama de implantação, por sua vez, mostra a distribuição física dos componentes de software sobre os nós de hardware, o que permite comunicar a arquitetura de execução do sistema \cite{omg2017uml}.

No S.R.I., a modelagem em UML foi utilizada para documentar o escopo funcional, a estrutura de dados e a dinâmica de interação do módulo de gerenciamento de locais, conforme apresentado no Capítulo \ref{cap:requisitos}.

\section{BPMN e a modelagem de processos de negócio}

A Business Process Model and Notation (BPMN) é uma notação gráfica padronizada para a representação de processos de negócio. Seu elemento central é o diagrama de processo, no qual as atividades são organizadas em raias que identificam o responsável por cada tarefa, permitindo representar de forma explícita a divisão de trabalho entre os participantes do processo \cite{omg2011bpmn}.

A notação distingue os eventos de início e de fim, as atividades, os gateways que representam decisões e os fluxos de sequência que conectam esses elementos. Um gateway exclusivo, por exemplo, indica que apenas um entre vários caminhos possíveis será seguido, conforme uma condição avaliada em tempo de execução \cite{omg2011bpmn}. A utilidade prática da BPMN em projetos de software está na possibilidade de documentar o processo na situação em que ele se encontra, convencionalmente chamada de \textit{as is}, e, a partir dessa descrição, identificar as atividades manuais, redundantes ou sujeitas a erro que justificam a automação.

No presente trabalho, a notação foi empregada para descrever o processo manual de coleta e entrega vigente na operação estudada, conforme detalhado no Capítulo \ref{cap:workflow-bpmn}.

\section{Desenvolvimento móvel multiplataforma com Flutter e Dart}

O desenvolvimento de aplicações móveis tradicionalmente exige a manutenção de bases de código distintas para cada sistema operacional, o que multiplica o esforço de implementação e de manutenção. Os frameworks multiplataforma surgem como alternativa a esse cenário, permitindo que um único código-fonte seja compilado para plataformas diferentes \cite{sommerville2019}.

O Flutter é um framework de interface de usuário que adota essa abordagem e compila o código para as plataformas Android e iOS a partir de uma mesma base \cite{flutter2024}. Sua arquitetura baseia-se em uma árvore de componentes, chamados \textit{widgets}, que descrevem tanto a estrutura quanto a aparência da interface, sendo recompostos sempre que o estado da aplicação se altera. Esse modelo de composição favorece a construção de telas consistentes e a reutilização de elementos visuais entre diferentes partes do aplicativo.

A linguagem utilizada pelo Flutter é o Dart, uma linguagem orientada a objetos, de tipagem estática e com suporte a programação assíncrona por meio de \textit{futures} e \textit{streams} \cite{dart2024}. Esse suporte é relevante para o S.R.I., pois as operações de consulta a serviços externos de geocodificação e de roteamento são inerentemente assíncronas e precisam ocorrer sem bloquear a interface enquanto o usuário interage com o mapa ou com a lista de locais.

\section{Persistência local com SQLite e sqflite}

A escolha do mecanismo de persistência em aplicações móveis precisa considerar a disponibilidade da conexão de rede, uma vez que o aplicativo pode ser utilizado em campo, em regiões com cobertura instável. Bancos de dados embarcados respondem a essa necessidade por armazenarem os dados no próprio dispositivo, dispensando um servidor remoto para as operações básicas \cite{sqlite2024}.

O SQLite é um mecanismo de banco de dados relacional embarcado, implementado como uma biblioteca e gravado em um único arquivo no sistema de arquivos do dispositivo. Por ser transacional e compatível com a linguagem SQL, oferece as garantias de integridade típicas de sistemas gerenciadores de banco de dados maiores, com custo de configuração praticamente nulo \cite{sqlite2024}. Em aplicações Flutter, o acesso ao SQLite é normalmente realizado por meio da biblioteca sqflite, que disponibiliza operações de criação e de migração de esquema, além das operações de inserção, consulta, atualização e exclusão de registros \cite{sqflite2024}.

Essa capacidade de execução de operações de escrita e de leitura localmente é o que permite que o cadastro de locais do S.R.I. funcione mesmo quando o dispositivo se encontra sem conexão, ficando a dependência de rede restrita às etapas de geocodificação e de cálculo de trajeto \cite{sriEspecComp2025}.

\section{Geolocalização, geocodificação e mapas}

A identificação da posição do usuário em tempo real é obtida por meio do Sistema Global de Posicionamento (GPS), presente na maioria dos dispositivos móveis atuais. O acesso a esses dados no Flutter é viabilizado por bibliotecas especializadas, que expõem a posição corrente do dispositivo e permitem o acompanhamento das mudanças de localização \cite{geolocator2024}. Além do posicionamento, o sistema precisa converter endereços textuais em coordenadas geográficas, operação denominada geocodificação, e o inverso, chamado geocodificação reversa.

O OpenStreetMap (OSM) é um projeto colaborativo que mantém uma base cartográfica mundial de uso livre, construída a partir de contribuições voluntárias e disponibilizada sob licença aberta \cite{osm2024}. Sobre essa base, o projeto mantém o Nominatim, um serviço de busca que realiza a geocodificação de endereços e de nomes de locais e devolve as coordenadas correspondentes \cite{nominatim2024}. A adoção desses serviços é vantajosa para o S.R.I. porque elimina a necessidade de contratação de chaves de API comerciais, alinhando-se à diretriz de custo mínimo declarada nos requisitos do sistema \cite{sriEspecComp2025}.

A renderização cartográfica no aplicativo é feita pela composição de imagens quadradas, denominadas \textit{tiles}, que o cliente solicita ao servidor de mapas conforme a área visível e o nível de aproximação. A biblioteca flutter\_map implementa esse mecanismo no Flutter, permitindo sobrepor marcadores e traçados sobre o mapa \cite{fluttermap2024}.

\section{Roteamento de veículos e heurísticas de solução}

A definição da sequência em que um conjunto de pontos deve ser visitado constitui um problema clássico de otimização combinatória. Na sua formulação mais simples, conhecida como problema do caixeiro-viajante (TSP, do inglês \textit{Traveling Salesman Problem}), busca-se o percurso de menor custo que visite cada ponto exatamente uma vez e retorne ao ponto de partida \cite{cormen2012}. Quando se considera uma frota de veículos com capacidade limitada, que deve atender a um conjunto de clientes a partir de um depósito, tem-se o problema de roteamento de veículos (VRP, do inglês \textit{Vehicle Routing Problem}), formulado originalmente por \citeonline{dantzig1959}.

Ambos os problemas pertencem à classe dos problemas NP-difíceis, de modo que a obtenção da solução ótima torna-se inviável, em tempo hábil, para instâncias de porte realista. Por essa razão, a prática da roteirização apoia-se em métodos heurísticos, que não garantem a optimalidade da solução, mas produzem resultados de qualidade aceitável em tempo compatível com a operação \cite{cunha2000}. Essa característica é especialmente relevante em aplicações móveis, nas quais o cálculo precisa ser concluído em poucos segundos e executado em hardware com recursos limitados.

Entre as heurísticas construtivas, o método do vizinho mais próximo parte de um ponto inicial e, a cada passo, seleciona como próximo destino o ponto ainda não visitado que se encontra à menor distância do ponto corrente. O método é simples de implementar e de custo computacional baixo, mas tende a produzir soluções de qualidade mediana, com frequência deixando para o final do percurso os pontos que se encontram mais distantes da origem \cite{cormen2012}. Uma variante imediata consiste em construir a rota de ida pelo método do vizinho mais próximo a partir da origem e, encerrada a sequência de paradas, acrescentar o retorno ao ponto inicial, o que produz o circuito completo utilizado nas operações de coleta.

Para melhorar soluções já construídas, empregam-se heurísticas de refinamento, entre as quais se destaca o método 2-opt, proposto por \citeonline{croes1958}. O procedimento remove duas arestas da rota e reconecta os trechos resultantes de modo inverso, verificando se o custo total do percurso diminui. Quando a troca representa uma melhoria, a nova configuração é mantida e o processo se repete até que nenhuma substituição vantajosa seja encontrada. A combinação de uma heurística construtiva com um método de refinamento é a estratégia usual para obter, em tempo aceitável, circuitos sensivelmente melhores do que aqueles produzidos pela construção isolada \cite{croes1958}.

\section{APIs de roteamento por vias: OSRM}

O cálculo da distância entre dois pontos pode ser aproximado pela distância em linha reta, mas essa medida ignora a malha viária e, portanto, não corresponde ao trajeto efetivamente percorrível. Para obter rotas aderentes às vias, utilizam-se serviços de roteamento que operam sobre uma base cartográfica e aplicam algoritmos de busca em grafo para determinar o caminho de menor custo entre dois pontos.

O Open Source Routing Machine (OSRM) é um motor de roteamento de código aberto que calcula rotas de veículos sobre os dados do OpenStreetMap, oferecendo uma interface de programação que devolve, para um conjunto ordenado de coordenadas, a geometria do trajeto, a distância e o tempo estimado de percurso \cite{osrm2024}. O desempenho do serviço decorre de um extenso pré-processamento da malha viária, que organiza a rede em estruturas hierárquicas e permite responder às consultas em tempo muito inferior ao exigido por algoritmos de busca aplicados diretamente ao grafo original \cite{luxen2011}.

Para o S.R.I., a integração com um serviço de roteamento por vias cumpre dupla função: fornece a distância real entre os pontos cadastrados, insumo necessário para a aplicação das heurísticas descritas na seção anterior, e devolve a geometria do percurso para a representação do trajeto sobre o mapa exibido ao usuário.
